\documentclass{article}

\begin{document}

\title{Spritz\_InitData}
\date{22/10/2015}
\maketitle

\abstract{Spritz\_InitData - the initial data for the Spritz code}

\section{Introduction}
\label{sec:intro}

This thorn generates the initial data for the Spritz code. 





\subsection{Magnetized Neutron Stars}
\label{sec:TOVMHD}

It can be useful sometimes to evolve magnetized neutron stars with a
magnetic field completely contained inside the star. Several
prescriptions are of course possible and here we have implemented the
most common one which is based on adding a purely poloidal field
having the following vector potential in cylindrical coordinates $(r,\phi,z)$:
%
\begin{eqnarray}
  A^r    &=& 0 \\
  A^\phi &=& A_b r^2 \mathrm{max}(P-P_{cut},0) \\
  A^z    &=& 0
\end{eqnarray}
%
where $A_b$ define the amplitude of the magnetic field and its value
is set through the parameter {\tt Ab}. $P_{cut}$ is equivalent to {\tt
press\_cut}$\times \mathrm{max}(P)$ and {\tt press\_cut} is a real
positive parameter which define the point inside the star where the
vector potential goes to zero.

Since Spritz works in cartesian coordinates it is necessary to express
$A^j$ in cartesian components obtaining the following expressions:

\begin{eqnarray}
  A^x &=& -r^3 A_b f(P) \mathrm{Sin}(\phi)\\
  A^y &=&  r^3 A_b f(P) \mathrm{Cos}(\phi)\\
  A^z &=& 0
\end{eqnarray}

where $r=\sqrt{(x-x_c)^2+(y-y_c)^2}$,
$\phi=\mathrm{atan}\left(\frac{y-y_c}{x-x_c}\right)$, $f(P)\equiv
\mathrm{max}(P-P_{cut},0)$, $x_c$ and $y_c$ are the coordinates of the
center of the star assumed to be located on the plane $z=0$. The
position of the center is automatically computed if the parameter {\tt
Spritz::compute\_DCandMaxima} is set equal to {\tt ``yes''} and it is
assumed to be coincident with the position of the maximum of the
rest-mass density $\rho$. Otherwise it is assumed to be located in the
origin of the grid.


The magnetic field components in cartesian coordinates are then simply
given by $\mathbf{B}=\frac{1}{\sqrt{\gamma}}\nabla \times \mathbf{A}$
or explicitly:

\begin{eqnarray}
B^x &=& -\frac{(x-x_c) A_b r^2}{\sqrt{\gamma}} \frac{\partial f(P)}{\partial z} \\
B^y &=& -\frac{(y-y_c) A_b r^2}{\sqrt{\gamma}} \frac{\partial f(P)}{\partial z} \\
B^z &=& \frac{A_b r^2}{\sqrt{\gamma}}\left[4f(P)+(y-y_c)\frac{\partial f(P)}{\partial y} +(x-x_c)\frac{\partial f(P)}{\partial x}\right]
\end{eqnarray}
%
where $\gamma$ is the determinant of the three-metric $\gamma_{ij}$.\\

In order to use this initial configuration one has to set:\\ 
%
{\tt Spritz\_InitData::Spritz\_initial\_data=''TOVMHD''}.\\
%

If used with Carpet then the parameter {\tt tovmhd\_timelevels} must
 be set equal to three.\\ 
%
{\bf NOTE: If used with Carpet you may see some small discontinuities
 (depending also on the resolution you use) in the magnetic field
 variables in the buffer zones of the finest grids (if the magnetic
 field crosses the boundary between refinement levels). This is due to
 the interpolators used for prolongation (ENO or WENO) which are not
 precise as the lagrangian interpolator.}

{\bf Unit Conversion: to convert the magnetic field from cactus units
to Gauss one needs simply to multiply by $8.34583 \times 10^{19}$.}
\begin{thebibliography}{1}

\bibitem{livrevsrrfd}
J.~M. Mart{\'{\i}} and E.~M{\"u}ller.
\newblock Numerical hydrodynamics in {S}pecial {R}elativity.
\newblock {\em Living Rev. Relativity}, {\bf 3}, 1999.
\newblock [Article in online journal], cited on 31/7/01,
  http://www.livingreviews.org/Articles/Volume2/1999-3marti/index.html.

\bibitem{hydro1}
J.~A. Font, M. Miller, W. Suen and M. Tobias.
\newblock Three Dimensional Numerical General Relativistic
Hydrodynamics I: Formulations, Methods, and Code Tests
\newblock {\em Phys. Rev.}, {\bf D61}, 044011, 2000.

\end{thebibliography}


\include{interface}
\include{param}
\include{schedule}

\end{document}
